\scanpage{198}

La partie a) dans c) suggère la question suivante :

e) Soit $A$ schéma abélien sur $S$, $\omega$ section de $A$, telle que pour tt pt fermé $s\in S$, $\omega(s)$ se remonte à $A_s^{(p)}$ via Verschiebung $A_s^{(\unsure{p/s|p/k|p/S})}\longrightarrow A_s$. Est-il vrai qu'alors $\omega$ est d'ordre fini ?

\deleted{Dans le cas où $S$ lui-\unsure{$=$|$\simeq$} est de car $p>0$ (au lieu de dominer $\mathrm{Spec}\,\mathbb Z$) on dispose de $A^{p/S}\longrightarrow A$, et l'image inverse de $\omega$ dans $A^{p/S}$ est un revêtt \unsure{de|ét}}

\margin{fin 1\textsuperscript{ere} partie,}

% archive sheet number: 197
